\newpage
\TFMChapter[even]{Estado del Arte y Fundamentos Teóricos.}

\TFMSection{Introducción}
El desarrollo de soluciones basadas en IA generativa requiere combinar modelos, técnicas y componentes tecnológicos. En el caso de este proyecto, la generación de propuestas de viaje se apoya en modelos de lenguaje, sistemas RAG, procesamiento inteligente de documentos, agentes especializados y servicios externos de información.

Este capítulo presenta los fundamentos teóricos necesarios para comprender la solución desarrollada y, finalmente, revisa algunas soluciones existentes relacionadas con la aplicación de la IA generativa en el sector turístico.

\TFMSection{Modelos de lenguaje y modelos generativos}
Los modelos de lenguaje de gran escala (\textit{Large Language Models}, LLM) constituyen el componente central de los sistemas de IA generativa actuales. Su arquitectura de referencia es el Transformer, introducido por \textcite{vaswani2017attention}, que sustituye las redes recurrentes por un mecanismo de atención (\textit{self-attention}) capaz de ponderar la relevancia relativa de todos los elementos de una secuencia de entrada de forma paralela. Esto permite entrenar modelos sobre volúmenes de texto mucho mayores y capturar dependencias de largo alcance de manera más eficiente que las arquitecturas recurrentes previas. A partir de esta arquitectura, los modelos generativos actuales operan de forma autorregresiva: dado un contexto de entrada, el modelo predice de forma iterativa el siguiente elemento de la secuencia (\textit{token}) hasta completar la respuesta, muestreando sobre una distribución de probabilidad condicionada al historial de tokens generados.

El procesamiento de texto por parte de estos modelos no se realiza a nivel de palabra completa, sino mediante una tokenización basada en subunidades léxicas. Se ha adoptado ampliamente la técnica de \textit{Byte Pair Encoding} (BPE), propuesta originalmente para traducción automática por \textcite{sennrich2016subword}, que permite representar un vocabulario abierto mediante un conjunto fijo de subpalabras, mitigando el problema de las palabras fuera de vocabulario. Esta segmentación en tokens define, a su vez, el concepto de ventana de contexto: el número máximo de tokens que el modelo puede procesar simultáneamente entre la entrada y la salida, una limitación arquitectónica relevante para el diseño de sistemas que combinan LLM con recuperación documental. Finalmente, el trabajo de \textcite{brown2020fewshot} sobre GPT-3 demostró que, al escalar el número de parámetros y el volumen de datos de entrenamiento, estos modelos desarrollan capacidades emergentes de aprendizaje en contexto (\textit{in-context learning}), pudiendo resolver tareas nuevas a partir de unos pocos ejemplos incluidos en el propio \textit{prompt} sin necesidad de reentrenamiento, propiedad determinante para las técnicas de adaptación sin ajuste de pesos descritas más adelante en este capítulo.

\TFMSection{Recuperación aumentada por generación}
La generación aumentada por recuperación (\textit{Retrieval-Augmented Generation}, RAG) fue formalizada por \textcite{lewis2020rag} como una arquitectura híbrida que combina un componente paramétrico (el modelo de lenguaje) con un componente no paramétrico (un índice de documentos externo), de manera que la generación de texto se condiciona a pasajes recuperados dinámicamente en lugar de depender exclusivamente del conocimiento almacenado en los pesos del modelo. Este enfoque permite fundamentar las respuestas en fuentes documentales verificables y actualizables, y mitiga --sin eliminar-- el problema de las alucinaciones inherente a los LLM.

El primer paso de cualquier sistema RAG es la representación semántica del contenido mediante \textit{embeddings}: vectores de alta dimensión que capturan el significado de un texto de forma que la proximidad geométrica entre vectores refleja similitud semántica. Para la codificación de pasajes y consultas se emplean habitualmente arquitecturas \textit{bi-encoder} como Sentence-BERT \parencite{reimers2019sentencebert}, que modifica BERT mediante una estructura siamesa para producir embeddings de frases comparables por similitud coseno de forma eficiente, y el enfoque de recuperación densa (\textit{Dense Passage Retrieval}, DPR) de \textcite{karpukhin2020dense}, que demuestra que un recuperador basado en embeddings puede superar de forma sustancial a los métodos léxicos clásicos en tareas de pregunta-respuesta de dominio abierto. Sobre estos embeddings se construye la búsqueda semántica, que sustituye la coincidencia exacta de términos por una búsqueda por vecinos más cercanos en el espacio vectorial.

\clearpage
Previamente a la indexación, los documentos deben someterse a un proceso de \textit{chunking} o fragmentación en unidades de tamaño manejable, de forma que cada fragmento contenga suficiente contexto semántico sin exceder los límites prácticos de recuperación y generación; no existe una estrategia universal, y la elección del tamaño y el criterio de segmentación condiciona directamente la calidad de la recuperación posterior \parencite{pinecone2023chunking}. El almacenamiento y la consulta eficiente de estos fragmentos vectorizados se apoya en bases de datos vectoriales, cuyo fundamento algorítmico procede de estructuras de búsqueda aproximada de vecinos más cercanos (\textit{Approximate Nearest Neighbor}, ANN): el índice de \textcite{johnson2019faiss} y los grafos jerárquicos de mundo pequeño navegable (HNSW) de \textcite{malkov2020hnsw} son las dos familias de algoritmos que sustentan la mayoría de motores vectoriales empleados en producción en la actualidad, entre ellos los servicios gestionados de recuperación documental como Amazon Bedrock \parencite{aws2026bedrockrag}. Tras la recuperación inicial de un conjunto amplio de candidatos, es habitual aplicar una etapa de \textit{reranking} que reordena los pasajes mediante un modelo más costoso basado en \textit{cross-encoding}, como el propuesto por \textcite{nogueira2019passage}, que procesa conjuntamente consulta y pasaje para obtener una puntuación de relevancia más precisa que la obtenida únicamente por similitud de embeddings. Finalmente, los pasajes seleccionados se incorporan al contexto del LLM para producir una generación de respuestas fundamentada, en la que el modelo debe sintetizar la información recuperada apoyándose explícitamente en las fuentes proporcionadas, tal y como se planteó originalmente en \textcite{lewis2020rag}.

\TFMSection{Sistemas basados en agentes}
Un agente inteligente basado en LLM se define como un sistema que, a diferencia de un modelo generativo puramente reactivo, es capaz de intercalar razonamiento y acción de forma iterativa para resolver tareas complejas. El paradigma de referencia es ReAct, propuesto por \textcite{yao2023react}, que muestra que generar explícitamente trazas de razonamiento (\textit{thoughts}) intercaladas con acciones ejecutables sobre el entorno mejora tanto la interpretabilidad como la fiabilidad del proceso de decisión frente a enfoques que separan razonamiento y actuación. Este patrón intercalado de pensar-actuar-observar constituye la base de la mayoría de arquitecturas de agentes actuales, incluida la empleada en este proyecto.

La capacidad de uso de herramientas es lo que permite a un agente superar las limitaciones de un LLM aislado (conocimiento desactualizado, incapacidad de ejecutar cálculos exactos o de acceder a sistemas externos). \textcite{schick2023toolformer} demostraron que un modelo de lenguaje puede aprender de forma autosupervisada cuándo y cómo invocar APIs externas e incorporar sus resultados a la generación posterior, sentando las bases conceptuales del \textit{tool calling} que hoy exponen de forma nativa las principales plataformas de modelos. En cuanto a la memoria y el estado, los agentes requieren mecanismos que permitan conservar información entre pasos de ejecución o entre intentos sucesivos ante un mismo objetivo; el trabajo de \textcite{shinn2023reflexion} introduce el concepto de refuerzo verbal, en el que el agente genera una autocrítica en lenguaje natural de sus intentos fallidos y la almacena en una memoria episódica para mejorar su desempeño posterior, sin necesidad de reentrenar los pesos del modelo. El enrutamiento y la toma de decisiones --esto es, decidir qué herramienta invocar, qué agente especializado debe intervenir o si la tarea está resuelta-- se apoya en estos mismos mecanismos de razonamiento intercalado descritos por \textcite{yao2023react}.

\clearpage
Cuando la complejidad de la tarea lo exige, resulta habitual descomponer el sistema en arquitecturas multiagente, en las que distintos agentes especializados --cada uno con su propio rol, herramientas y contexto-- colaboran mediante intercambio de mensajes o mediante traspaso de control. El marco AutoGen \parencite{wu2023autogen} formaliza este patrón como una conversación entre agentes configurables que integran LLM, herramientas y, opcionalmente, intervención humana. Para orquestar estos flujos de manera explícita y controlable, en lugar de dejar la coordinación implícita en el propio texto conversacional, este proyecto emplea LangGraph \parencite{langchain2026langgraph}, un framework que representa el sistema como un grafo de estados donde los nodos son agentes o funciones y las aristas codifican las transiciones posibles, lo que permite modelar tanto flujos secuenciales como bucles condicionales, puntos de control humano y persistencia del estado a lo largo de una interacción de larga duración.

\TFMSection{Procesamiento inteligente de documentos}
El procesamiento de la documentación de una agencia de viajes (tarifarios, catálogos de actividades, fichas de alojamiento) presenta un reto adicional respecto al texto plano: gran parte de la información relevante está codificada en el diseño visual del documento (tablas, columnas, cabeceras) y no únicamente en su contenido textual. El enfoque clásico para abordar este problema combina un motor de reconocimiento óptico de caracteres (OCR) con un modelo de comprensión de \textit{layout}; \textcite{xu2020layoutlm} muestran con LayoutLM que preentrenar conjuntamente sobre texto y posición espacial de cada palabra en la página mejora sustancialmente la extracción de información en documentos escaneados frente a modelos que solo consideran la secuencia textual. Una alternativa más reciente, y cada vez más adoptada, elimina por completo la dependencia de un motor OCR intermedio: \textcite{kim2022donut} proponen Donut, un Transformer que procesa directamente la imagen del documento de extremo a extremo, evitando así la propagación de errores de reconocimiento de caracteres hacia las fases posteriores del pipeline.

Este mismo principio --trabajar directamente sobre la representación visual del documento en lugar de sobre texto ya extraído-- es el que subyace a los modelos multimodales de propósito general actuales, capaces de aceptar imágenes o páginas de PDF convertidas a imagen como entrada nativa junto con instrucciones en lenguaje natural. Tanto el informe técnico de GPT-4 \parencite{openai2023gpt4} como la documentación oficial de capacidades de visión de modelos comparables \parencite{anthropic2026vision} describen esta capacidad multimodal, que en este proyecto se emplea para convertir cada página de un documento PDF en una imagen y solicitar al modelo la extracción estructurada de su contenido (texto, tablas, metadatos) en un único paso, sustituyendo así el pipeline tradicional de OCR más reglas de post-procesado. La normalización del contenido extraído a un formato intermedio como Markdown, así como la extracción de metadatos estructurados, constituyen decisiones de ingeniería del pipeline sobre las que no existe un cuerpo de literatura académica específico; se trata de una práctica consolidada en la industria para uniformizar la salida de modelos multimodales heterogéneos antes de la fase de preparación para RAG, en la que el documento normalizado se somete al proceso de fragmentación y vectorización descrito en la sección anterior \parencite{lewis2020rag,pinecone2023chunking}.

\TFMSection{Técnicas de especialización de modelos mediante ingeniería de prompts}
Frente al reentrenamiento o ajuste fino de los pesos de un modelo, existe un conjunto de técnicas que permiten especializar el comportamiento de un LLM preentrenado exclusivamente mediante el diseño de su entrada, sin modificar sus parámetros. \textcite{liu2023pretrain} sistematizan este paradigma bajo el término \textit{prompt-based learning}, estableciendo una taxonomía de las estrategias de diseño de instrucciones que condicionan la tarea que el modelo debe resolver. Dentro de esta familia de técnicas, el prompt de sistema constituye el mecanismo mediante el cual se fija de forma persistente el rol, las restricciones y el estilo de respuesta del modelo para toda una sesión, independientemente de las instrucciones del usuario, y es el mecanismo empleado en este proyecto para acotar el comportamiento de cada agente especializado.

El \textit{few-shot prompting}, cuya eficacia fue documentada empíricamente por \textcite{brown2020fewshot} al observar que GPT-3 podía resolver tareas nuevas a partir de unos pocos ejemplos incluidos en el propio prompt, permite orientar el formato y el criterio de salida del modelo sin necesidad de ejemplos de entrenamiento adicionales. De forma complementaria, técnicas como el razonamiento en cadena (\textit{chain-of-thought}) propuesto por \textcite{wei2022chainofthought} muestran que solicitar al modelo que explicite sus pasos intermedios de razonamiento antes de emitir una respuesta final mejora su desempeño en tareas que requieren varios pasos de inferencia. El uso de herramientas --descrito en la sección anterior a partir del trabajo de \textcite{schick2023toolformer}-- y la recuperación aumentada por generación \parencite{lewis2020rag} constituyen, en este mismo sentido, formas de adaptación en tiempo de inferencia: en ambos casos se amplía el conocimiento o las capacidades del modelo inyectando información o resultados externos en su contexto, en lugar de modificar sus pesos.

\clearpage
Para que la salida de un LLM pueda integrarse de forma fiable en un sistema software --por ejemplo, para generar un documento de propuesta comercial con una estructura predefinida-- resulta necesario forzar salidas estructuradas que respeten un esquema determinado; la documentación oficial sobre uso de herramientas y salidas estructuradas de los proveedores de modelos \parencite{anthropic2026toolsuse} describe los mecanismos mediante los cuales se restringe la generación del modelo a un formato válido conforme a un esquema, lo que permite además una primera capa de validación sintáctica de la respuesta antes de cualquier verificación semántica adicional. Esta validación es especialmente relevante dado que los LLM pueden producir contenido no fundamentado en las fuentes proporcionadas (alucinaciones), fenómeno caracterizado de forma extensa por \textcite{ji2023hallucination}, por lo que la validación de respuestas --ya sea estructural, mediante esquemas, o semántica, contrastando la salida con los documentos recuperados-- constituye una capa de control necesaria en cualquier sistema de producción. Por último, la gestión del contexto --qué información se incluye, en qué orden y hasta qué límite de tokens-- resulta crítica en sistemas que combinan histórico de conversación, resultados de herramientas y pasajes recuperados dentro de una misma ventana de contexto finita \parencite{vaswani2017attention}.

\TFMSection{Tecnologías para la integración de fuentes externas}
La conexión de un LLM con sistemas y datos externos de forma estandarizada es un problema que, hasta hace poco, cada aplicación resolvía mediante integraciones \textit{ad hoc}. El Protocolo de Contexto de Modelo (\textit{Model Context Protocol}, MCP), presentado por Anthropic en 2024 \parencite{anthropic2024mcpannounce} y formalizado en su especificación técnica \parencite{mcp2025specification}, define un estándar abierto de comunicación cliente-servidor mediante el cual una aplicación que aloja un modelo de lenguaje puede descubrir y utilizar de forma uniforme herramientas, fuentes de datos y prompts expuestos por servidores externos, sustituyendo así las integraciones propietarias fragmentadas por un único protocolo interoperable.

Entre estos servicios externos se encuentran las APIs de proveedores de contenido turístico, como las expuestas por Amadeus for Developers \parencite{amadeus2026developers}, que proporcionan acceso mediante API REST a la búsqueda de ofertas de vuelos y de disponibilidad y precios hoteleros sobre el inventario de un \textit{Global Distribution System}. La integración de este tipo de APIs como herramientas invocables por los agentes --siguiendo el patrón de uso de herramientas descrito anteriormente \parencite{schick2023toolformer,yao2023react}-- permite fundamentar la generación de propuestas comerciales en disponibilidad y precios reales y actualizados, en lugar de depender del conocimiento estático del modelo de lenguaje.

\TFMSection{Marco legal, ético y de seguridad}
El tratamiento de documentación de clientes por parte de un sistema de IA generativa en el sector turístico se encuentra sujeto, en primer lugar, al Reglamento General de Protección de Datos \parencite{gdpr2016}, que establece los principios de licitud, minimización de datos, limitación de la finalidad y las bases jurídicas exigibles para el tratamiento de datos personales que un sistema de este tipo maneja de forma inherente, así como las obligaciones de transparencia frente al interesado y las garantías exigibles cuando el tratamiento se apoya en proveedores de modelos de lenguaje operados por terceros. Junto a la protección de datos personales, se plantea la cuestión de la propiedad y confidencialidad de la documentación de la agencia y de sus clientes cuando esta se envía a servicios de inferencia externos, lo que exige políticas claras de retención, cifrado y control de qué proveedor puede utilizar dichos datos para fines distintos de la propia inferencia.

Un segundo bloque de riesgos deriva de las limitaciones intrínsecas de los LLM, en particular la generación de contenido plausible pero no fundamentado en las fuentes disponibles, fenómeno descrito de forma sistemática por \textcite{ji2023hallucination}. En un sistema que genera propuestas comerciales con precios, fechas y disponibilidad, una alucinación no detectada tiene consecuencias contractuales directas, lo que justifica las capas de validación de respuestas y de fundamentación mediante RAG descritas en las secciones anteriores \parencite{lewis2020rag}. Esto se conecta directamente con el principio de transparencia: el usuario final debe poder distinguir el contenido generado automáticamente del contenido verificado por un agente humano, y disponer de trazabilidad sobre qué fuentes documentales han fundamentado cada afirmación de la propuesta.

\clearpage
Por último, la seguridad de credenciales y el control de accesos resultan especialmente relevantes en una arquitectura que integra múltiples servicios externos mediante protocolos como MCP \parencite{mcp2025specification,anthropic2024mcpannounce}, en la que cada servidor de herramientas actúa como intermediario con capacidad de ejecutar acciones en sistemas de terceros; esto exige aplicar el principio de mínimo privilegio en la concesión de credenciales a cada agente y validar el origen y el alcance de cada invocación de herramienta. En el ámbito europeo, este conjunto de riesgos se enmarca además en el Reglamento de Inteligencia Artificial \parencite{euai2024act}, que clasifica los sistemas de IA por niveles de riesgo y establece obligaciones de transparencia, documentación técnica y supervisión humana proporcionales al riesgo del sistema; si bien un asistente de generación de propuestas de viaje no se encuadra, en principio, entre los sistemas de alto riesgo enumerados en dicho Reglamento, sí le resultan aplicables las obligaciones de transparencia hacia el usuario cuando este interactúa con un sistema de IA generativa.

\TFMSection{Estado del Arte y Soluciones Existentes}
La aplicación de la IA generativa al sector turístico ha sido objeto de un interés creciente tanto desde la industria como desde la investigación académica. Desde la perspectiva académica, \textcite{dogru2025genai} presentan, a partir de un panel amplio de investigadores del ámbito de \textit{hospitality} y turismo, un marco conceptual que identifica las principales líneas de aplicación de la IA generativa en el sector (personalización de la experiencia del viajero, generación de contenido, automatización de tareas operativas) junto con los retos abiertos de investigación. En una línea más aplicada y cercana al problema concreto de este TFM --la generación asistida de itinerarios y propuestas de viaje--, \textcite{volchek2024chatgpt} evalúan de forma empírica la capacidad de ChatGPT para actuar como planificador de itinerarios, comparando los resultados generados por el modelo con itinerarios elaborados por expertos turísticos; sus resultados muestran que el contenido generado es accesible y fácil de entender, pero significativamente menos preciso y específico que el de un experto humano, lo que refuerza la necesidad de fundamentar la generación mediante fuentes verificadas en lugar de confiar únicamente en el conocimiento paramétrico del modelo. En la misma línea, \textcite{ilieva2024genai} evalúan con profesionales del sector y viajeros el efecto de las herramientas de IA generativa sobre la eficiencia operativa de las agencias y la personalización de la experiencia del cliente, confirmando beneficios claros en la generación de contenido, pero señalando también limitaciones de actualización en tiempo real que motivan arquitecturas híbridas de modelo de lenguaje más recuperación documental.

Desde la perspectiva de la evaluación de capacidades, \textcite{xie2024travelplanner} introducen TravelPlanner, un banco de pruebas específicamente diseñado para medir la capacidad de agentes basados en LLM para resolver tareas de planificación de viajes con restricciones realistas de presupuesto, preferencias y disponibilidad; sus resultados iniciales --una tasa de éxito de apenas el 0,6\,\% para GPT-4 actuando como agente autónomo sin arquitectura de apoyo-- evidencian que la planificación de viajes de extremo a extremo mediante un LLM sin una arquitectura de agentes, herramientas y verificación adecuada dista de ser un problema resuelto, lo que refuerza la pertinencia de un enfoque multiagente con orquestación explícita, como el adoptado en este proyecto, frente a un enfoque de prompt único.

\clearpage
Desde la perspectiva de la industria, dos de los mayores operadores de distribución turística en línea han desplegado públicamente asistentes conversacionales de generación de propuestas de viaje basados en LLM: Expedia lanzó en 2023 una experiencia de planificación conversacional integrada en su aplicación y, posteriormente, en el propio ChatGPT \parencite{expedia2023chatgpt}; y Booking.com lanzó de forma paralela su \textit{AI Trip Planner}, que combina modelos de aprendizaje automático propios con tecnología de modelos de lenguaje para ofrecer una experiencia conversacional de planificación integrada directamente en el flujo de reserva de alojamiento \parencite{booking2023openai}. Ambos casos constituyen evidencia real y de gran escala de que el patrón arquitectónico ``modelo de lenguaje conversacional más recuperación de inventario propio más flujo de reserva'' --análogo, salvando las distancias de escala, al enfoque de este TFM-- ha sido validado por actores relevantes del sector como una dirección viable, si bien ninguno de los dos casos publica detalles técnicos de su arquitectura interna más allá de los anuncios corporativos aquí referenciados, por lo que no ha sido posible verificar si emplean patrones de agentes multipaso, RAG documental o arquitecturas de grafos comparables a las descritas en este capítulo.
